\section{The cyclic feedback map}
For $n\ge1$, put $X_n=\{1,2,3\}^n$ and read position subscripts modulo
$n$. Define $R:X_n\to X_n$ by letting $R(x)_i$ count the strict
left-to-right records of
\[
 x_i,x_{i+1},\ldots,x_{i+n-1}.
\]
The first letter is a record, and a later letter is a record only when it
exceeds every preceding letter. Thus $1\le R(x)_i\le3$. All scans use the
same old word. Positions are labelled: words differing by a rotation are
not identified. We seek the exact images and recurrent dynamics of $R$,
and all sources of a prescribed one-step target.

Strict and weak record statistics for words and multisets, including
single-scan record counts and position templates, are classical
\citep{myers2008maxima}. Only the first occurrence of a value can be a
strict record. Repeatedly jumping to the next strictly greater value
therefore traces a scan's records; this uses the strict nearest-value
primitive studied in \citet{berkman1993nearest}, after reversing the order.
Neither the statistic nor its static pointer description is new here.
The feedback question recomputes these counts at every cyclic start.
For example, $123123$ and $123213$ have the same multiset, first-occurrence
order and displayed single-scan record positions, but their cyclic outputs
are $321321$ and $321221$. Those conditioning data alone do not determine
the joint target; this observation does not exclude more elaborate reductions.

We prove an exact two-step reflection core and, by a separate source-maximum
decomposition, an evaluated one-step inverse for each target. Transfer-matrix
counts, weak binary chain counts and the classical integer-product optimum
are elementary consequences and tools, not separate contributions.
